\documentclass[11pt,a4paper]{article}
\usepackage[utf8]{inputenc}
\usepackage[margin=1in]{geometry}
\usepackage{amsmath,amssymb,booktabs,graphicx,hyperref,natbib}

\title{\textbf{Hydroclimatic Anomalies, First-Mile Supply-Chain Bottlenecks, and Crop Resilience: Multi-Scale Evidence from 26-Year Satellite Reanalysis, 40,538 Infrastructure Works, and Debiased Machine Learning in Peru}}
\author{\textbf{Juan Marcos Vilchez-Canchari}$^{a,b,*}$ \\
\small $^a$ \textit{Escuela de Posgrado, Universidad C\'esar Vallejo, Lima, Peru} \\
\small $^b$ \textit{Universidad Privada San Juan Bautista, Lima, Peru} \\
\small $^*$ \textit{Corresponding author. Email address: jvilchezc@ucvvirtual.edu.pe}}
\date{November 2026}

\begin{document}
\maketitle

\begin{abstract}
Agro-export supply chains and smallholder food systems across the Global South face escalating compound risks from hydroclimatic anomalies and chronic public infrastructure procurement failures. This study develops a multi-scale empirical and operations-research framework for Peru by integrating 624 monthly ECMWF Copernicus ERA5-HEAT satellite reanalysis grids (Universal Thermal Climate Index [UTCI] downscaled with district topographic lapse rates and nocturnal Tropical Nights $>20^\circ$C from 2000 to 2025), 312 monthly BCRP/MIDAGRI agricultural output series across 24 series (3 macro-sectoral GDP indices, 3 aggregate baskets, and 18 individual export and staple crops), and 40,538 geo-referenced public investment projects (17,351 hydro-agricultural irrigation works and 23,187 rural transport supply-chain works) across all 1,873 districts. Five-fold cross-fitted Double/Debiased Machine Learning (DML) conditioning on district altitude, topographic slope, log population, rurality, and department fixed effects demonstrates that infrastructure exposures vary across specifications, whereas direct cross-sectional variance in irrigation bottlenecks (DML theta = +0.098, p = 0.6971; OLS-FE beta = +0.049, p = 0.8543), first-mile logistics bottlenecks (DML theta = -0.063, p = 0.7995; OLS-FE beta = -0.288, p = 0.2753), intra-district ERA5 hydroclimatic shocks (DML theta = -0.052, p = 0.8929; OLS-FE beta = +0.250, p = 0.4500), compound climate x bottleneck interaction (DML theta = -0.086, p = 0.7639; OLS-FE beta = -0.124, p = 0.6694) is largely absorbed by topographic and departmental controls, and completed irrigation and feeder-road buffer density is associated with DML theta = -0.328 (p = 0.2089; OLS-FE beta = -0.573, p = 0.0346). Under a S/ 2,500 million completion budget across 1,497 bottlenecked districts, our DML-weighted Pareto optimization captures 57.8\% of national resilience gains, outperforming historical inertia (21.9\%) by +35.9 percentage points.
\end{abstract}

\section{Introduction}
See complete manuscript in \texttt{submission\_package\_Q1\_Elsevier/04\_Main\_Manuscript.docx} and tables/figures in \texttt{paper/tables/} and \texttt{paper/figures/}.

\end{document}
